% GENERATED FILE. Edit canonical lessons, then run npm run generate:books.
% canonical-source-hash: fnv1a64:7f40c896e441cd81
% canonical-lessons: ZH-C102-yinliao, ZH-C102-guozhi, ZH-C102-kele, ZH-C102-pijiu, ZH-C102-binggan

\chapter{Something to Drink}
\label{ch:zh-something-to-drink}
\hlchaptermodality{102}

\begin{chapteropening}
\textbf{By the end of this chapter:} \emph{I can ask for a drink: juice, cola or beer, and a biscuit.}

Complete the last lesson of chapter 102: I can ask for a drink: juice, cola or beer, and a biscuit.
\end{chapteropening}

\section[yǐnliào]{yǐnliào — drink; soft drink}

\label{lesson:ZH-C102-yinliao}



\begin{quote}
Before the new one: say the Mandarin for plate, then the Mandarin for knife.
\end{quote}

\subsection*{You'll want to know: yǐnliào}
\emph{yǐnliào} — \textquotedblleft{}drink; soft drink\textquotedblright{}.

\begin{quote}
\emph{nǐ xiǎng hē shénme yǐnliào} — What would you like to drink?
\end{quote}

Its characters come later; for now it is a word for the ear.

\begin{cousinweb}
\emph{yǐnliào} is built from yǐn 'to drink' and liào, both inherited from Middle Chinese, \emph{'imX} and \emph{lew} (Guangyun, in Baxter's transcription). The dictionary glosses the second 'to measure' and lists it with a departing tone too, \emph{lewH}, the one liào continues; how the pair came to mean a drink is not settled.
\end{cousinweb}

\begin{grammarlens}[title={What you've built}]
One more word for asking for what I want.
\end{grammarlens}

\subsection*{Your turn}
\begin{itemize}
\raggedright
  \item \emph{Say it:} \emph{yǐnliào}
  \item \emph{Say it:} \emph{yǐnliào}, once more
  \item \emph{Say it:} \emph{nǐ xiǎng hē shénme yǐnliào} — What would you like to drink?
  \item \emph{Recall:} say the Mandarin for plate, then the Mandarin for knife, then say \emph{yǐnliào} again
\end{itemize}

\begin{tcolorbox}[breakable,colback=teal!4,colframe=teal!35!black,title={Before you move on}]
What does yǐnliào mean? (\textquotedblleft{}Drink; soft drink\textquotedblright{}.) Say it once more.
\end{tcolorbox}
\section[guǒzhī]{guǒzhī — fruit juice}

\label{lesson:ZH-C102-guozhi}



\begin{quote}
Before the new one: say the Mandarin for knife, then the Mandarin for soft drink.
\end{quote}

\subsection*{You'll want to know: guǒzhī}
\emph{guǒzhī} — \textquotedblleft{}fruit juice\textquotedblright{}.

\begin{quote}
\emph{wǒ yào yì bēi guǒzhī} — A glass of juice for me, please.
\end{quote}

Its characters come later; for now it is a word for the ear.

\begin{cousinweb}
\emph{guǒzhī} is built from guǒ 'fruit' and zhī 'juice', both inherited from Middle Chinese, \emph{kwaX} and \emph{tsyip} (Guangyun, in Baxter's transcription), glossed 'the fruit of a tree' and 'liquid'.
\end{cousinweb}

\begin{grammarlens}[title={What you've built}]
One more word for asking for what I want.
\end{grammarlens}

\subsection*{Your turn}
\begin{itemize}
\raggedright
  \item \emph{Say it:} \emph{guǒzhī}
  \item \emph{Say it:} \emph{guǒzhī}, once more
  \item \emph{Say it:} \emph{wǒ yào yì bēi guǒzhī} — A glass of juice for me, please.
  \item \emph{Recall:} say the Mandarin for knife, then the Mandarin for soft drink, then say \emph{guǒzhī} again
\end{itemize}

\begin{tcolorbox}[breakable,colback=teal!4,colframe=teal!35!black,title={Before you move on}]
What does guǒzhī mean? (\textquotedblleft{}Fruit juice\textquotedblright{}.) Say it once more.
\end{tcolorbox}
\section[kělè]{kělè — cola}

\label{lesson:ZH-C102-kele}



\begin{quote}
Before the new one: say the Mandarin for soft drink, then the Mandarin for juice.
\end{quote}

\subsection*{You'll want to know: kělè}
\emph{kělè} — \textquotedblleft{}cola\textquotedblright{}.

\begin{quote}
\emph{tāmen dōu hē kělè} — They all drink cola.
\end{quote}

Its characters come later; for now it is a word for the ear.

\begin{cousinweb}
\emph{kělè} is borrowed from the name of a cola drink, usually traced to English \emph{cola}; its two syllables are written with the characters for kě 'may' and lè 'joy', chosen for their sound.
\end{cousinweb}

\begin{grammarlens}[title={What you've built}]
One more word for asking for what I want.
\end{grammarlens}

\subsection*{Your turn}
\begin{itemize}
\raggedright
  \item \emph{Say it:} \emph{kělè}
  \item \emph{Say it:} \emph{kělè}, once more
  \item \emph{Say it:} \emph{tāmen dōu hē kělè} — They all drink cola.
  \item \emph{Recall:} say the Mandarin for soft drink, then the Mandarin for juice, then say \emph{kělè} again
\end{itemize}

\begin{tcolorbox}[breakable,colback=teal!4,colframe=teal!35!black,title={Before you move on}]
What does kělè mean? (\textquotedblleft{}Cola\textquotedblright{}.) Say it once more.
\end{tcolorbox}
\section[píjiǔ]{píjiǔ — beer}

\label{lesson:ZH-C102-pijiu}



\begin{quote}
Before the new one: say the Mandarin for juice, then the Mandarin for cola.
\end{quote}

\subsection*{You'll want to know: píjiǔ}
\emph{píjiǔ} — \textquotedblleft{}beer\textquotedblright{}.

\begin{quote}
\emph{wǒ bù hē píjiǔ} — I don't drink beer.
\end{quote}

Its characters come later; for now it is a word for the ear.

\begin{cousinweb}
\emph{píjiǔ} is built from pí and jiǔ 'alcoholic drink'. The second is inherited from Middle Chinese \emph{tsjuwX} (Guangyun, in Baxter's transcription); the first is borrowed for its sound from a European word for beer, usually traced to English \emph{beer}.
\end{cousinweb}

\begin{grammarlens}[title={What you've built}]
One more word for asking for what I want.
\end{grammarlens}

\subsection*{Your turn}
\begin{itemize}
\raggedright
  \item \emph{Say it:} \emph{píjiǔ}
  \item \emph{Say it:} \emph{píjiǔ}, once more
  \item \emph{Say it:} \emph{wǒ bù hē píjiǔ} — I don't drink beer.
  \item \emph{Recall:} say the Mandarin for juice, then the Mandarin for cola, then say \emph{píjiǔ} again
\end{itemize}

\begin{tcolorbox}[breakable,colback=teal!4,colframe=teal!35!black,title={Before you move on}]
What does píjiǔ mean? (\textquotedblleft{}Beer\textquotedblright{}.) Say it once more.
\end{tcolorbox}
\section[bǐnggān]{bǐnggān — biscuit; cracker}

\label{lesson:ZH-C102-binggan}



\begin{quote}
Before the new one: say the Mandarin for cola, then the Mandarin for beer.
\end{quote}

\subsection*{You'll want to know: bǐnggān}
\emph{bǐnggān} — \textquotedblleft{}biscuit; cracker\textquotedblright{}.

\begin{quote}
\emph{wǒ xiǎngyào yí kuài bǐnggān} — I'd like a biscuit.
\end{quote}

Its characters come later; for now it is a word for the ear.

\begin{cousinweb}
\emph{bǐnggān} is built from bǐng, a flat cake, and gān 'dry', both inherited from Middle Chinese, \emph{pjiengX} and \emph{kan} (Guangyun, in Baxter's transcription); gān is the gān of gānjìng 'clean'.
\end{cousinweb}

\begin{grammarlens}[title={What you've built}]
One more word for asking for what I want.
\end{grammarlens}

\subsection*{Your turn}
\begin{itemize}
\raggedright
  \item \emph{Say it:} \emph{bǐnggān}
  \item \emph{Say it:} \emph{bǐnggān}, once more
  \item \emph{Say it:} \emph{wǒ xiǎngyào yí kuài bǐnggān} — I'd like a biscuit.
  \item \emph{Recall:} say the Mandarin for cola, then the Mandarin for beer, then say \emph{bǐnggān} again
\end{itemize}

\begin{tcolorbox}[breakable,colback=teal!4,colframe=teal!35!black,title={Before you move on}]
What does bǐnggān mean? (\textquotedblleft{}Biscuit; cracker\textquotedblright{}.) Say it once more.
\end{tcolorbox}
